\chapter{Magnetic Forces, Vector Potentials, and the Action Principle}

\lectureoverview{Electrostatic forces fit the familiar conservative pattern, but magnetic forces depend on velocity. To place them inside Lagrangian mechanics, we first develop the needed vector calculus, use the source-free character of the magnetic field to introduce a vector potential, examine its gauge freedom, and then construct the magnetic term in the action.}

\section{A New Kind of Force}

The preceding discussion of angular momentum, gyroscopes, and Poisson brackets brought several themes together. In particular, a rotating charged body in a magnetic field precesses in much the same way as a gyroscope. We now turn from that analogy to the motion of an individual charged particle.

Electrostatic forces require no new mechanics. If \(U(\mathbf x)\) is a potential energy, then
\begin{equation}
  \mathbf F=-\nabla U(\mathbf x),
\end{equation}
and the usual argument gives conservation of kinetic plus potential energy. For a particle of charge \(e\) in an electrostatic potential \(V(\mathbf x)\),
\begin{equation}
  U(\mathbf x)=eV(\mathbf x),
  \qquad
  \mathbf F_{\mathrm{el}}=-e\nabla V.
\end{equation}
Magnetic forces have a different form. They depend on the particle's velocity, so they cannot be introduced merely by adding another ordinary potential energy. Nevertheless, they do follow from an action principle and admit both Lagrangian and Hamiltonian descriptions. Reaching that description requires a short but essential interlude in vector calculus.

\section{Fields and Local Operations}

\subsection{What a field records}

A field assigns a physical quantity to every point of space and, in general, to every time. Temperature and pressure are scalar fields. Wind velocity is a vector field. Electric and magnetic fields are vector fields. Other theories contain scalar fields even at microscopic scales, while gravitation is described by a tensor field. For the present discussion we suppress any time dependence and write fields as functions of \(\mathbf x\).

The operational point of a field is locality. Imagine that a certain region is hidden inside a bag. To measure the electric field in that region, put a small test charge there and observe its acceleration. The force is determined by the field where the charge is located, not by a separate consultation with every distant source. The same statement applies to a magnetic field and a moving test charge. Later, the vector potential will force us to qualify this simple picture, but it is the right place to begin.

\subsection{The derivative operator}

Derivatives turn fields into other fields. We collect the three spatial derivatives into
\begin{equation}
  \nabla
  =
  \left(
    \frac{\partial}{\partial x},
    \frac{\partial}{\partial y},
    \frac{\partial}{\partial z}
  \right).
\end{equation}
The symbol is read \emph{del} or \emph{nabla}. Its upside-down-triangle appearance is memorable, but its real significance is that it behaves formally like a vector whose components are differential operators.

This is ordinary three-dimensional physical space, not phase space and not a many-particle configuration space. Moreover, \(\nabla\) is not a vector field that can stand alone. It must act on something. If \(T(\mathbf x)\) is a scalar temperature field, then
\begin{equation}
  \nabla T
  =
  \left(
    \frac{\partial T}{\partial x},
    \frac{\partial T}{\partial y},
    \frac{\partial T}{\partial z}
  \right)
\end{equation}
is its gradient, a genuine vector field.

\begin{figure}[htbp]
\centering
\includegraphics[width=0.82\textwidth]{lecture_09_figure_01.png}
\caption{The board introduces \(\nabla\) as a vector of spatial derivative operators and begins applying it to a scalar field \(T(\mathbf x)\).}
\lectureframe{00:11:30}
\end{figure}

A scalar is not merely a quantity written with one component. It is a quantity whose value is independent of the orientation of the coordinate axes. Temperature is the standard example. Given a vector field \(\mathbf E\), the local expression \(\mathbf E\cdot\mathbf E\) is also a scalar. By contrast, constructing a potential by integrating a field is a nonlocal operation; it uses values along a path.

\subsection{Divergence}

For a vector field \(\mathbf A=(A_x,A_y,A_z)\), the scalar
\begin{equation}
  \nabla\cdot\mathbf A
  =
  \frac{\partial A_x}{\partial x}
  +\frac{\partial A_y}{\partial y}
  +\frac{\partial A_z}{\partial z}
\end{equation}
is its divergence. In fluid language, positive divergence measures a local tendency to flow outward, while negative divergence measures a tendency to converge inward. A source of fluid produces positive divergence; a sink produces negative divergence. This picture is qualitative, but it will make the source-free character of the magnetic field transparent.

\section{Cross Products, Epsilon Symbols, and Curl}

\subsection{The cross product in components}

For two vectors \(\mathbf V\) and \(\mathbf A\), one component of the cross product is
\begin{equation}
  (\mathbf V\times\mathbf A)_x
  =
  V_yA_z-V_zA_y.
\end{equation}
The other two follow by cycling \(x\to y\to z\to x\). It is useful to replace \(x,y,z\) by the indices \(1,2,3\). The Kronecker symbol is
\begin{equation}
  \delta_{ij}
  =
  \begin{cases}
    1,&i=j,\\
    0,&i\ne j,
  \end{cases}
\end{equation}
and the Levi--Civita symbol is defined by
\begin{equation}
  \epsilon_{123}=\epsilon_{231}=\epsilon_{312}=1,
  \qquad
  \epsilon_{132}=\epsilon_{321}=\epsilon_{213}=-1,
\end{equation}
with \(\epsilon_{ijk}=0\) whenever any two indices coincide. A cyclic ordering around a three-hour clock gives the plus signs; reversing the direction gives the minus signs. Interchanging any pair of indices reverses the sign, so \(\epsilon_{ijk}\) is completely antisymmetric.

We also use the summation convention: an index repeated in a term is summed from \(1\) to \(3\). The three cross-product equations then become
\begin{equation}
  (\mathbf V\times\mathbf A)_i
  =
  \epsilon_{ijk}V_jA_k.
\end{equation}
For example, setting \(i=1\) leaves only \((j,k)=(2,3)\) and \((3,2)\):
\begin{equation}
  (\mathbf V\times\mathbf A)_1
  =
  \epsilon_{123}V_2A_3
  +\epsilon_{132}V_3A_2
  =
  V_2A_3-V_3A_2.
\end{equation}
The compact notation therefore contains exactly the familiar component rule.

\subsection{A vector disguised as an antisymmetric matrix}

There is a second representation:
\begin{equation}
  V_iA_j-V_jA_i
  =
  \epsilon_{ijk}(\mathbf V\times\mathbf A)_k.
\end{equation}
The left side has two indices and is antisymmetric under \(i\leftrightarrow j\). As a matrix it has the pattern
\begin{equation}
  M_{ij}
  =
  \begin{pmatrix}
    0&M_{12}&M_{13}\\
    -M_{12}&0&M_{23}\\
    -M_{13}&-M_{23}&0
  \end{pmatrix}.
\end{equation}
The diagonal entries vanish, and the lower triangle is fixed by the upper triangle. There are therefore three independent components.

\begin{figure}[htbp]
\centering
\includegraphics[width=0.82\textwidth]{lecture_09_figure_02.png}
\caption{One cross-product component and the schematic \(3\times3\) antisymmetric array used to count its independent entries. The marks in the cells are mnemonic rather than additional symbols.}
\lectureframe{00:34:20}
\end{figure}

\begin{classroomqa}
\audiencequestion{Does the cross product, or the identification with another vector, make sense outside three dimensions?}
\lecturerresponse{The antisymmetric two-index object makes sense in any dimension, but its identification with a vector is special to three dimensions. A vector in three dimensions has three components, and an antisymmetric \(3\times3\) matrix has three independent components. In four dimensions the corresponding counts are four and six, since \(4(4-1)/2=6\). The antisymmetric object still exists and is naturally regarded as a two-form, but it is not another four-vector.}
\end{classroomqa}

\subsection{Curl}

Replacing \(\mathbf V\) by \(\nabla\) defines the curl:
\begin{equation}
  (\nabla\times\mathbf A)_i
  =
  \epsilon_{ijk}\frac{\partial A_k}{\partial x_j}.
\end{equation}
In ordinary components,
\begin{align}
  (\nabla\times\mathbf A)_x
  &=
  \frac{\partial A_z}{\partial y}
  -\frac{\partial A_y}{\partial z},\\
  (\nabla\times\mathbf A)_y
  &=
  \frac{\partial A_x}{\partial z}
  -\frac{\partial A_z}{\partial x},\\
  (\nabla\times\mathbf A)_z
  &=
  \frac{\partial A_y}{\partial x}
  -\frac{\partial A_x}{\partial y}.
\end{align}
Where divergence measures local spreading, curl measures local circulation. A fluid rotating around an axis has nonzero curl along that axis. The picture is particularly vivid for velocity fields, though the definition applies to any vector field.

\section{Two Structural Identities}

Two facts will do nearly all of the mathematical work:
\begin{equation}
  \nabla\cdot(\nabla\times\mathbf A)=0,
  \qquad
  \nabla\times(\nabla s)=0.
\end{equation}
The first follows because the second derivatives are symmetric while \(\epsilon_{ijk}\) is antisymmetric:
\begin{equation}
  \nabla\cdot(\nabla\times\mathbf A)
  =
  \epsilon_{ijk}\partial_i\partial_j A_k
  =0.
\end{equation}
The second follows from equality of mixed partial derivatives:
\begin{equation}
  \bigl[\nabla\times(\nabla s)\bigr]_i
  =
  \epsilon_{ijk}\partial_j\partial_k s
  =0.
\end{equation}

\begin{figure}[htbp]
\centering
\includegraphics[width=0.80\textwidth]{lecture_09_figure_03.png}
\caption{The first identity on the board: the divergence of every curl vanishes.}
\lectureframe{00:38:26}
\end{figure}

These are the easy directions, suitable for the lecture's promised homework or pop quiz. The converse statements require more care. Locally, or globally on an appropriate simply connected or contractible domain,
\begin{align}
  \nabla\cdot\mathbf F=0
  &\quad\Longrightarrow\quad
  \mathbf F=\nabla\times\mathbf A,\\
  \nabla\times\mathbf F=0
  &\quad\Longrightarrow\quad
  \mathbf F=\nabla s.
\end{align}
Without the domain qualification, holes and other global topological features can obstruct these conclusions. The lecture uses the local form, which is exactly what is needed below.

\section{Why the Magnetic Field Is a Curl}

\subsection{No observed magnetic sources}

A magnetic field \(\mathbf B\) is defined operationally through its effect on charged particles and currents. In ordinary electromagnetism without magnetic monopoles,
\begin{equation}
  \nabla\cdot\mathbf B=0.
\end{equation}
This differs from the electric field, whose divergence is proportional to electric charge density. Electric charges are sources and sinks of electric flux. No isolated source or sink of magnetic flux has been observed.

Magnetic monopoles would change this equation. Unified field theories often provide theoretical reasons to expect such objects, but if they exist they may be too massive to produce in available experiments. The present theory therefore proceeds with zero magnetic charge density.

\begin{classroomqa}
\audiencequestion{What would a magnetic monopole be? Ordinary magnets already have north and south poles.}
\lecturerresponse{An isolated monopole would be a genuine source or sink of magnetic field, analogous to an electric charge. An ordinary bar magnet is different. Field lines that emerge near its north pole curve through space and return through the south pole, continuing through the material of the magnet. What appears externally to be a source and a sink is connected, like fluid moving through a hidden pipe. There is no net creation or destruction of magnetic flux.}
\end{classroomqa}

\begin{figure}[htbp]
\centering
\includegraphics[width=0.76\textwidth]{lecture_09_figure_04.png}
\caption{The board contrasts apparent north and south sources while joining their field lines through the magnet, emphasizing that the magnetic flux does not begin or end there.}
\lectureframe{00:44:56}
\end{figure}

Since \(\mathbf B\) is divergence-free, the converse theorem lets us write, at least locally,
\begin{equation}
  \boxed{\mathbf B=\nabla\times\mathbf A}.
\end{equation}
The field \(\mathbf A\) is the vector potential. At first this may look like an unnecessary re-description of the measurable field \(\mathbf B\). Its purpose will become clear only when we construct the action.

\subsection{The vector potential is not unique}

If \(\mathbf A\) is one vector potential, then
\begin{equation}
  \mathbf A'=\mathbf A+\nabla s
\end{equation}
is another, because
\begin{equation}
  \nabla\times\mathbf A'
  =
  \nabla\times\mathbf A+\nabla\times(\nabla s)
  =
  \mathbf B.
\end{equation}
This replacement is a \emph{gauge transformation}. The magnetic field is gauge invariant.

\begin{figure}[htbp]
\centering
\includegraphics[width=0.80\textwidth]{lecture_09_figure_05.png}
\caption{The defining relation \(\mathbf B=\nabla\times\mathbf A\) and the gauge shift \(\mathbf A'=\mathbf A+\nabla s\) appear together on the board.}
\lectureframe{01:03:01}
\end{figure}

It is pronounced \emph{gauge}, not \emph{gouache} or \emph{gouge}--a distinction the lecture enforces with mock severity. The word itself is historical. It arose from an early attempt to connect the freedom with comparisons of length standards at different locations; that original interpretation did not survive, but the name did.

\begin{classroomqa}
\audiencequestion{In this discussion, does a gauge transformation mean specifically \(\mathbf A\mapsto\mathbf A+\nabla s\)?}
\lecturerresponse{Yes. That is the complete electromagnetic transformation being used here. Gauge transformations in Yang--Mills theory and other field theories generalize the same idea, but electromagnetism needs only the addition of a scalar gradient.}
\end{classroomqa}

The nonuniqueness also clarifies what is directly measurable. A charged particle in a small region responds to the local \(\mathbf B\), and all vector potentials related by a gauge transformation predict the same force there. There is no local experiment of this classical kind that assigns a unique value to \(\mathbf A\).

\begin{classroomqa}
\audiencequestion{Is adding a gradient the complete description of the nonuniqueness?}
\lecturerresponse{Locally on a simple domain, yes. If \(\nabla\times\mathbf A'=\nabla\times\mathbf A\), then \(\nabla\times(\mathbf A'-\mathbf A)=0\), so the difference is a gradient. Globally, a multiply connected domain can carry additional subtleties. Adding a constant vector is already included, since a constant vector is the gradient of a linear scalar function. Adding a constant to \(s\) itself changes nothing because its gradient vanishes.}
\end{classroomqa}

\begin{classroomqa}
\audiencequestion{If gauge freedom has an underlying physical reason, what is it?}
\lecturerresponse{The lecture postpones that deeper interpretation. The useful first step is to work with the freedom and see how it enters the action. Once the formalism is familiar, the reason gauge descriptions are indispensable becomes easier to recognize.}
\end{classroomqa}

\section{One Magnetic Field, Many Potentials}

Consider the simplest possible example,
\begin{equation}
  \mathbf B=b\,\hat{\mathbf z},
  \qquad
  B_x=B_y=0,
  \qquad
  B_z=b,
\end{equation}
where \(b\) is constant. One vector potential is
\begin{equation}
  \mathbf A^{(1)}=(0,bx,0).
\end{equation}
Indeed,
\begin{equation}
  (\nabla\times\mathbf A^{(1)})_z
  =
  \frac{\partial A_y^{(1)}}{\partial x}
  -\frac{\partial A_x^{(1)}}{\partial y}
  =b,
\end{equation}
and the other curl components vanish.

This choice seems to single out the \(y\)-direction even though the physical field is rotationally symmetric in the \(x\)-\(y\) plane. A second choice makes the arbitrariness explicit:
\begin{equation}
  \mathbf A^{(2)}=(-by,0,0).
\end{equation}
It gives
\begin{equation}
  (\nabla\times\mathbf A^{(2)})_z
  =
  0-\frac{\partial(-by)}{\partial y}
  =b.
\end{equation}
The two potentials differ by
\begin{equation}
  \mathbf A^{(2)}-\mathbf A^{(1)}
  =
  (-by,-bx,0)
  =
  \nabla(-bxy).
\end{equation}
They are therefore gauge equivalent. A rotationally symmetric choice,
\begin{equation}
  \mathbf A^{(\mathrm{sym})}
  =
  \frac12\,\mathbf B\times\mathbf r
  =
  \left(-\frac{by}{2},\frac{bx}{2},0\right),
\end{equation}
is another member of the same family.

\begin{figure}[htbp]
\centering
\includegraphics[width=0.80\textwidth]{lecture_09_figure_06.png}
\caption{Two board panels hold the choices \(A_y=bx\) and \(A_x=-by\), both producing the same uniform field \(B_z=b\).}
\lectureframe{01:07:20}
\end{figure}

\begin{classroomqa}
\audiencequestion{Is this the same ambiguity as arranging the current-carrying coils in different ways to produce a field?}
\lecturerresponse{No. Different arrangements of current are different physical sources and may produce different fields elsewhere, even if they agree in one region. Gauge freedom remains after the magnetic field has been specified everywhere: the same physical \(\mathbf B\) still has many vector potentials.}
\end{classroomqa}

\begin{classroomqa}
\audiencequestion{Can one picture adding a gradient as stretching or compressing without adding a twist?}
\lecturerresponse{That is a useful geometric intuition. Curl detects the local twisting part, and a pure gradient has zero curl. The statement is encoded exactly in \(\nabla\times\nabla s=0\).}
\end{classroomqa}

\begin{classroomqa}
\audiencequestion{Is there a fixed relation between the direction of \(\mathbf A\) and the direction of \(\mathbf B\)?}
\lecturerresponse{No gauge-independent relation exists. The examples above happen to have \(A_z=0\), but a constant \(A_z\) could be added without changing any derivative. Then \(\mathbf A\) would no longer be perpendicular to \(\mathbf B\). Only the curl has invariant physical content in this classical discussion.}
\end{classroomqa}

There is a milder analogue in electrostatics. Since \(\mathbf E=-\nabla V\), replacing \(V\) by \(V+C\) does not change \(\mathbf E\). The term \emph{vector potential} was chosen partly by analogy with this freedom in the ordinary scalar potential: one may change the potential in a way that leaves the physical field unchanged.

\section{The Lorentz Force}

We can now state the force law that the action principle must reproduce. In the conventions used here,
\begin{equation}
  \boxed{
  \mathbf F
  =
  e\left(
    \mathbf E+\frac{\mathbf v}{c}\times\mathbf B
  \right)
  =
  -e\nabla V+\frac{e}{c}\mathbf v\times\mathbf B
  }.
\end{equation}
The symbol \(e\) is the particle's electric charge, \(\mathbf v=\dot{\mathbf x}\) is its velocity, and \(V\) is the scalar potential. The notation must not tempt us to confuse \(V\) with \(\mathbf v\). In SI language \(V\) is measured in volts and \(e\) in coulombs; the displayed \(1/c\) belongs to the unit convention used in the lecture. The structural point is independent of units.

\begin{figure}[htbp]
\centering
\includegraphics[width=0.82\textwidth]{lecture_09_figure_07.png}
\caption{The electric and magnetic terms of the Lorentz force are displayed together. The magnetic contribution is proportional to \(\mathbf v\times\mathbf B\).}
\lectureframe{00:58:54}
\end{figure}

The electric term is conservative and already familiar. The magnetic term is new. It vanishes for a stationary particle and is perpendicular to both \(\mathbf v\) and \(\mathbf B\). Consequently,
\begin{equation}
  \mathbf v\cdot\mathbf F_{\mathrm{mag}}
  =
  \frac{e}{c}\mathbf v\cdot(\mathbf v\times\mathbf B)
  =0.
\end{equation}
A static magnetic field bends a trajectory but does no work on the particle. This is unlike viscous drag, which is directed opposite the velocity and removes kinetic energy.

The factor \(v/c\) also indicates why, for comparable field scales in these conventions, magnetic effects are often smaller than electric effects for slowly moving particles. Electrons are light enough to reach large speeds, so magnetic deflection remains easy to observe.

Historically, forces on current-carrying wires were known before the force on an individual charge was isolated. Lorentz recognized that the wire force could be understood by assigning the velocity-dependent term to each moving charge. The complete electric-plus-magnetic expression is now called the Lorentz force law, although the magnetic piece is the novel part in the present discussion.

\begin{classroomqa}
\audiencequestion{Does the whole expression have a separate symbol, or is only the magnetic piece the Lorentz force?}
\lecturerresponse{The whole expression is simply the force \(\mathbf F=m\mathbf a\), and it is commonly called the Lorentz force. The magnetic term is the part historically associated most directly with Lorentz. Electric and magnetic fields are deeply related by relativity, but that unification is not needed for the present mechanical problem.}
\end{classroomqa}

\subsection{The Coriolis analogy}

We have seen a velocity crossed with a fixed vector before. In a frame rotating with angular velocity \(\boldsymbol\omega\), the Coriolis force can be written
\begin{equation}
  \mathbf F_{\mathrm{Cor}}
  =
  2m\,\mathbf v\times\boldsymbol\omega.
\end{equation}
It has the same algebraic form as the magnetic force:
\begin{equation}
  2m\boldsymbol\omega
  \quad\longleftrightarrow\quad
  \frac{e}{c}\mathbf B.
\end{equation}
A uniform magnetic field perpendicular to a plane therefore bends charged-particle trajectories in much the same way that a rotating frame deflects moving bodies.

\begin{figure}[htbp]
\centering
\includegraphics[width=0.82\textwidth]{lecture_09_figure_08.png}
\caption{The Lorentz force remains on the left while the board adds the corresponding Coriolis form \(2m\,\mathbf v\times\boldsymbol\omega\).}
\lectureframe{01:02:08}
\end{figure}

\begin{classroomqa}
\audiencequestion{The Coriolis force is called fictitious. Is it less real than the magnetic force because it depends on velocity?}
\lecturerresponse{Velocity dependence is not the issue. A person on a rotating platform can hold a moving object and feel the Coriolis force. It is called fictitious because it disappears in an inertial frame, while the rotating frame is defined precisely by the presence of such inertial terms. The word \emph{real} is often less useful than stating which frame and which equation are being used.}
\end{classroomqa}

The analogy can be pushed into a small thought experiment. On a rotating platform, a suitable electrostatic field can cancel the centrifugal term for a chosen particle, while a suitable uniform magnetic field can cancel the Coriolis term. The cancellation works only for a chosen charge-to-mass ratio: electric and magnetic forces scale with \(e\), whereas the inertial forces scale with \(m\). A field arrangement tuned for electrons would not cancel the same terms for particles with a different \(e/m\).

\section{Constructing the Charged-Particle Action}

\subsection{Stationary action with fixed endpoints}

We now ask the central question: can the Lorentz force be derived from an action principle? The answer is yes, but the vector potential is indispensable. A Lagrangian built locally from the particle and \(\mathbf B\) alone does not produce the required magnetic force. Rather than prove that negative statement, we construct the successful action directly.

For a trajectory \(\mathbf x(t)\), write
\begin{equation}
  S[\mathbf x]
  =
  \int_{t_i}^{t_f}L(\mathbf x,\dot{\mathbf x},t)\,dt.
\end{equation}
The lecture uses the same handwritten letter for the action and the vector potential, but we use \(S\) for the action and \(\mathbf A\) for the vector potential. The physical path makes \(S\) stationary under variations that hold both endpoints fixed. The traditional phrase \emph{least action} is convenient, though stationary action is more accurate because the extremum need not be a minimum.

The fixed endpoints matter. One may picture the ends of the path nailed to the board while the interior is wiggled. This is a trajectory in ordinary mechanics, not a string-theory object; the string image is only a way to visualize the variation.

\subsection{What may enter the Lagrangian?}

With all fields turned off, the Lagrangian is the kinetic energy,
\begin{equation}
  L_0=\frac{m}{2}\dot{\mathbf x}^{\,2}.
\end{equation}
The electrostatic term is also standard:
\begin{equation}
  L_{\mathrm{el}}=-eV(\mathbf x).
\end{equation}
The magnetic term must satisfy several clues:
\begin{enumerate}
  \item it must be proportional to the charge \(e\);
  \item it must contain the factor \(1/c\) in the lecture's units;
  \item it must involve the velocity;
  \item it must be a rotational scalar.
\end{enumerate}

The tempting trial \(\mathbf B\cdot\dot{\mathbf x}\) is a scalar and contains a velocity, but it gives the wrong equation of motion. The correct term uses the vector potential:
\begin{equation}
  L_{\mathrm{mag}}
  =
  \frac{e}{c}\mathbf A(\mathbf x)\cdot\dot{\mathbf x}
  =
  \frac{e}{c}A_i(\mathbf x)\dot x_i.
\end{equation}
Thus the full time-independent electromagnetic Lagrangian is
\begin{equation}
  \boxed{
  L
  =
  \frac{m}{2}\dot{\mathbf x}^{\,2}
  +\frac{e}{c}\mathbf A(\mathbf x)\cdot\dot{\mathbf x}
  -eV(\mathbf x)
  }.
\end{equation}
Lecture 10 will apply the Euler--Lagrange equations and show explicitly that this expression produces the Lorentz force.

\begin{figure}[htbp]
\centering
\includegraphics[width=0.82\textwidth]{lecture_09_figure_09.png}
\caption{The completed board construction of the action: a free kinetic term plus \((e/c)A_i\dot x_i\).}
\lectureframe{01:26:36}
\end{figure}

\section{Gauge Transformations Change Only the Endpoints}

The magnetic contribution to the action is
\begin{equation}
  S_{\mathrm{mag}}
  =
  \frac{e}{c}
  \int_{t_i}^{t_f}A_i(\mathbf x)\dot x_i\,dt.
\end{equation}
Since \(\dot x_i\,dt=dx_i\), it is a line integral along the particle's path \(\gamma\):
\begin{equation}
  S_{\mathrm{mag}}
  =
  \frac{e}{c}\int_\gamma A_i\,dx_i
  =
  \frac{e}{c}\int_\gamma\mathbf A\cdot d\mathbf x.
\end{equation}
Break the path into short segments, take the component of \(\mathbf A\) along each segment, multiply by the segment length, and add. That is the geometric meaning of the integral.

Now perform the gauge transformation
\begin{equation}
  A_i\longmapsto A_i+\partial_i s.
\end{equation}
The action changes by
\begin{align}
  \Delta S
  &=
  \frac{e}{c}\int_\gamma \partial_i s\,dx_i\\
  &=
  \frac{e}{c}\int_\gamma ds\\
  &=
  \frac{e}{c}
  \left[
    s(\mathbf x_f)-s(\mathbf x_i)
  \right].
\end{align}
Each infinitesimal segment contributes the value of \(s\) at one end minus its value at the other. On summing, all interior contributions cancel telescopically, leaving only the two endpoints.

\begin{figure}[htbp]
\centering
\includegraphics[width=0.82\textwidth]{lecture_09_figure_10.png}
\caption{The line integral, its division into small path segments, and the endpoint result \(s(\mathbf x_f)-s(\mathbf x_i)\).}
\lectureframe{01:32:31}
\end{figure}

\begin{classroomqa}
\audiencequestion{If the action changes under a gauge transformation, why does the trajectory not change?}
\lecturerresponse{The action principle compares paths with fixed initial and final points. The added quantity depends only on those endpoints, so it has the same value for every path in the comparison. Its variation is zero. The numerical action is not strictly invariant, but the Euler--Lagrange equations and the physical trajectory are gauge invariant.}
\end{classroomqa}

This distinction is the main conceptual result of the lecture:
\begin{equation}
  \begin{aligned}
    S&\ \hbox{may shift by an endpoint term},\\
    \delta S&\ \hbox{and the equations of motion do not}.
  \end{aligned}
\end{equation}
Gauge freedom is therefore compatible with the action principle. The vector potential can enter the Lagrangian even though no unique \(\mathbf A\) is selected by the magnetic field.

\section{Where the Argument Continues}

The present lecture has assembled the machinery and the action, but it deliberately stops before varying that action. The next steps are now fixed:
\begin{enumerate}
  \item apply the Euler--Lagrange equations and recover
  \(\mathbf F_{\mathrm{mag}}=(e/c)\mathbf v\times\mathbf B\);
  \item compute the canonical momentum, which will differ from \(m\mathbf v\);
  \item construct the Hamiltonian;
  \item identify the conservation laws;
  \item solve the motion of a charged particle in a uniform magnetic field.
\end{enumerate}

The route to those results explains why the vector potential was introduced. The magnetic field is the locally measurable curl, but the vector potential is the object that couples naturally to a particle's path. Gauge transformations alter that coupling only at the fixed endpoints, leaving the mechanics unchanged.
